\documentclass[aspectratio=169]{beamer}
\usepackage[english]{babel}
\usepackage[utf8]{inputenc}
\usepackage[T1]{fontenc}
\usepackage{lmodern}
\usepackage{listings}
\usepackage{xcolor}
\usepackage{graphicx}
\usepackage{textcomp}
\DeclareUnicodeCharacter{2014}{---}
\DeclareUnicodeCharacter{2013}{--}
\DeclareUnicodeCharacter{2212}{-}
\DeclareUnicodeCharacter{2192}{\ensuremath{\rightarrow}}
\DeclareUnicodeCharacter{00B1}{\ensuremath{\pm}}
\DeclareUnicodeCharacter{00D7}{\ensuremath{\times}}
\DeclareUnicodeCharacter{00B7}{\ensuremath{\cdot}}
\DeclareUnicodeCharacter{20AC}{EUR}
\DeclareUnicodeCharacter{2070}{\textsuperscript{0}}
\DeclareUnicodeCharacter{00B9}{\textsuperscript{1}}
\DeclareUnicodeCharacter{00B2}{\textsuperscript{2}}
\DeclareUnicodeCharacter{00B3}{\textsuperscript{3}}
\DeclareUnicodeCharacter{2074}{\textsuperscript{4}}
\DeclareUnicodeCharacter{2075}{\textsuperscript{5}}
\DeclareUnicodeCharacter{2076}{\textsuperscript{6}}
\DeclareUnicodeCharacter{2077}{\textsuperscript{7}}
\DeclareUnicodeCharacter{2078}{\textsuperscript{8}}
\DeclareUnicodeCharacter{2079}{\textsuperscript{9}}
\DeclareUnicodeCharacter{2248}{\ensuremath{\approx}}
\DeclareUnicodeCharacter{2264}{\ensuremath{\le}}

% Colors for code
\definecolor{codebg}{RGB}{245,245,245}
\definecolor{codecomment}{RGB}{0,128,0}
\definecolor{codekeyword}{RGB}{0,0,255}
\definecolor{codestring}{RGB}{163,21,21}
\definecolor{bandcolor}{RGB}{200,50,50}

\lstdefinestyle{pythonstyle}{
    language=Python,
    backgroundcolor=\color{codebg},
    basicstyle=\ttfamily\scriptsize,
    keywordstyle=\color{codekeyword},
    commentstyle=\color{codecomment},
    stringstyle=\color{codestring},
    showstringspaces=false,
    breaklines=true,
    frame=single,
    rulecolor=\color{black},
    escapeinside={(*@}{@*)},
    moredelim=**[l][\color{bandcolor}\bfseries]{\#\ ---\ NEW},
    moredelim=**[l][\color{bandcolor}\bfseries]{\#\ ---\ CHANGED},
    literate=
      {á}{{\'a}}1 {é}{{\'e}}1 {í}{{\'i}}1 {ó}{{\'o}}1 {ú}{{\'u}}1
      {Á}{{\'A}}1 {É}{{\'E}}1 {Í}{{\'I}}1 {Ó}{{\'O}}1 {Ú}{{\'U}}1
      {ñ}{{\~n}}1 {Ñ}{{\~N}}1 {ü}{{\"u}}1 {Ü}{{\"U}}1
      {¿}{{?`}}1 {¡}{{!`}}1
      {—}{{---}}1 {–}{{--}}1 {−}{{-}}1
      {±}{{\ensuremath{\pm}}}1 {×}{{\ensuremath{\times}}}1
      {→}{{\ensuremath{\rightarrow}}}1
      {°}{{\textdegree}}1
      {·}{{\ensuremath{\cdot}}}1
      {€}{{\texteuro}}1
      {≈}{{\ensuremath{\approx}}}1
      {≤}{{\ensuremath{\le}}}1
      {⁰}{{\textsuperscript{0}}}1 {¹}{{\textsuperscript{1}}}1
      {²}{{\textsuperscript{2}}}1 {³}{{\textsuperscript{3}}}1
      {⁴}{{\textsuperscript{4}}}1 {⁵}{{\textsuperscript{5}}}1
      {⁶}{{\textsuperscript{6}}}1 {⁷}{{\textsuperscript{7}}}1
      {⁸}{{\textsuperscript{8}}}1 {⁹}{{\textsuperscript{9}}}1
}

\lstset{style=pythonstyle}

% Title slide macro: \lessontitle{Lesson Number}{Title}
\newcommand{\lessontitle}[2]{
    \title{Lesson #1: #2}
    \author{}
    \institute{Tec de Monterrey}
    \begin{frame}[plain]
        \titlepage
    \end{frame}
}

% Figure frame macro: \figureframe{Title}{Image path}
\newcommand{\figureframe}[2]{
    \begin{frame}{#1}
        \begin{center}
            \includegraphics[height=0.8\textheight,width=\textwidth,keepaspectratio]{#2}
        \end{center}
    \end{frame}
}


% Listings pulled from src/ with \lstinputlisting, so the slide is the file.
\lstset{
    basicstyle=\ttfamily\fontsize{8pt}{9.6pt}\selectfont,
    breaklines=true,
    breakatwhitespace=true,
    columns=fullflexible,
    keepspaces=true,
    xleftmargin=0.3em,
    framesep=2pt,
    aboveskip=0pt,
    belowskip=0pt
}

\newcommand{\resultframe}[3]{%
    \begin{frame}{#1}
        \begin{center}
            \includegraphics[height=0.62\textheight,width=\textwidth,keepaspectratio]{#2}
        \end{center}
        \vspace{0.25em}
        {\small #3\par}
    \end{frame}%
}

\newcommand{\claimframe}[2]{%
    \begin{frame}{#1}
        {\small #2\par}
    \end{frame}%
}

\date{}

\begin{document}

\lessontitle{06}{Effectiveness}

\section{This lesson}

\begin{frame}{This lesson}
Lesson 03 postponed run-level evidence for selection pressure. This lesson is that instrument: success rate, uncertainty, budget.

\vspace{0.6em}
\begin{itemize}
    \item A success rate needs a verdict from outside the run, and a target with volume.
    \item One plug-in standard error is tens of points at eight runs and about one at a thousand; it is not a confidence interval.
    \item Effectiveness and efficiency pull opposite ways.
\end{itemize}
\end{frame}

% ================= success_rate_01_one_run =================
\begin{frame}{A rate is an estimate}
For \(s\) successes in \(n\) independent runs,
\[
\hat p=\frac{s}{n},\qquad
\widehat{\operatorname{SE}}(\hat p)=
\sqrt{\frac{\hat p(1-\hat p)}{n}}.
\]
This is one estimated standard deviation, not a confidence interval. The lesson
also reports a 95\% Wilson interval, which still expresses uncertainty at
0/\(n\) or \(n\)/\(n\). The grid optimum is a sampled reference.
\end{frame}

\begin{frame}{The 95\% Wilson interval}
With (z=1.96), let
\[
d=1+z^2/n,\quad c=\frac{\hat p+z^2/(2n)}d,\quad
h=\frac z d\sqrt{\frac{\hat p(1-\hat p)}n+\frac{z^2}{4n^2}}.
\]
The interval is ([c-h,c+h]). Unlike the plug-in SE alone, it remains
non-degenerate at 0 successes or all successes.
\end{frame}

\section{One run, and what it is worth}

\begin{frame}{One run, and what it is worth}
Every lesson so far ended by reading a single run. Lesson 01 read one run and
concluded the algorithm works; Lesson 02 read one run and concluded a stopping
rule is free; Lesson 03 and Lesson 05 both closed with an open question they
could not answer, and both questions had the same shape: does this change help,
over runs? This lesson is the machinery for answering that, and it starts by
showing why the question cannot be dodged.
\end{frame}

\resultframe{One run, and what it is worth}{../figures/success_rate_01_one_run.png}{%
    Spread across seeds is 0.8946 — 1.5× the tracked run's whole gain (+0.5822); seed 3 reports -0.0129, seed 4 reports -0.9076, and both reports are true}

% ================= success_rate_02_the_verdict =================
\section{What counts as success}

\begin{frame}{What counts as success}
Step 1 showed eight answers to one question. Before they can be averaged into
anything, one word has to be defined: success. A run cannot judge itself - it
reports the best thing it found, never how good that was - so the judgement has
to come from outside, and Lesson 02 step 5 built exactly that instrument: brute
force the landscape on a grid, and measure every run against the result. It also
left a warning attached to it, in its own closing lines: that reference "works
because x is one number; from Lesson 08 onwards it is not, and the honest answer
to 'did this run succeed?' becomes genuinely hard to get".
\end{frame}

\begin{frame}[fragile]{(1) brute\_force\_optimum()}
\lstinputlisting[basicstyle=\ttfamily\fontsize{7pt}{8.5pt}\selectfont,firstline=178,lastline=196]{../src/success_rate_02_the_verdict.py}
\end{frame}

\begin{frame}[fragile]{(2) TOLERANCE}
\lstinputlisting[basicstyle=\ttfamily\fontsize{8pt}{9.6pt}\selectfont,firstline=199,lastline=210]{../src/success_rate_02_the_verdict.py}
\end{frame}

\begin{frame}[fragile]{(3) within\_tolerance\_share()}
\lstinputlisting[basicstyle=\ttfamily\fontsize{6pt}{7.2pt}\selectfont,firstline=213,lastline=247]{../src/success_rate_02_the_verdict.py}
\end{frame}

\begin{frame}[fragile]{(4) SINE}
\lstinputlisting[basicstyle=\ttfamily\fontsize{8pt}{9.6pt}\selectfont,firstline=86,lastline=101]{../src/success_rate_02_the_verdict.py}
\end{frame}

\resultframe{What counts as success}{../figures/success_rate_02_the_verdict.png}{%
    The 2-D optimum is real (+0.500000) but occupies 0.003197\% of the grid — all 128 winning points sit exactly on y = 10 — so 0 of 8 runs "succeed"; on sine 1-D the target is 1.4299\% of the box (0.286 wide) and 5 of 8 runs find it: 62.5\% ± 17.1\%}

% ================= success_rate_03_a_thousand_runs =================
\section{A thousand runs}

\begin{frame}{A thousand runs}
Step 2 ended with a success rate measured on eight runs - 62.5\% - and a
warning about what such a rate is worth: if 62.5\% were the true rate, the
count in 8 runs would still wander by 1.4 runs from one experiment to the
next. This step pays that warning off. It runs the experiment a thousand
times, the way Gridin's Chapter 6 does, and replaces the anecdote with a
measurement: a rate, its standard error, and the full distribution of
answers behind it.
\end{frame}

\begin{frame}[fragile]{(1) success\_rate()}
\lstinputlisting[basicstyle=\ttfamily\fontsize{8pt}{9.6pt}\selectfont,firstline=185,lastline=198]{../src/success_rate_03_a_thousand_runs.py}
\end{frame}

\begin{frame}[fragile]{(2) standard\_error()}
\lstinputlisting[basicstyle=\ttfamily\fontsize{8pt}{9.6pt}\selectfont,firstline=201,lastline=211]{../src/success_rate_03_a_thousand_runs.py}
\end{frame}

\begin{frame}[fragile]{(3) RUNS = 1000}
\lstinputlisting[basicstyle=\ttfamily\fontsize{8pt}{9.6pt}\selectfont,firstline=53,lastline=53]{../src/success_rate_03_a_thousand_runs.py}
\end{frame}

\resultframe{A thousand runs}{../figures/success_rate_03_a_thousand_runs.png}{%
    The observed rate is 79.4\%; one plug-in SE is 1.3\% and the script reports a separate 95\% Wilson interval. The 8-run estimate differs by 16.9 points; 71/1000 runs end on Lesson 01's local peak at x = -4.51}

% ================= success_rate_04_the_budget =================
\section{The budget}

\begin{frame}{The budget}
Step 3 measured effectiveness: a success rate of 79.4\% at population 10, and
it closed with the question a rate cannot answer - what did those evaluations
cost, and were they well spent? This step counts them. Gridin's Chapter 6
does it with a class attribute on Individual that counts every fitness
evaluation, and with a sweep of population sizes that plots the final fitness
against the number of individuals evaluated.
\end{frame}

\begin{frame}[fragile]{(1) the evaluation counter}
\lstinputlisting[basicstyle=\ttfamily\fontsize{7pt}{8.5pt}\selectfont,firstline=77,lastline=98]{../src/success_rate_04_the_budget.py}
\end{frame}

\begin{frame}[fragile]{(2) the population sweep}
\lstinputlisting[basicstyle=\ttfamily\fontsize{7pt}{8.5pt}\selectfont,firstline=211,lastline=235]{../src/success_rate_04_the_budget.py}
\end{frame}

\begin{frame}[fragile]{(3) blind\_success()}
\lstinputlisting[basicstyle=\ttfamily\fontsize{8pt}{9.6pt}\selectfont,firstline=193,lastline=203]{../src/success_rate_04_the_budget.py}
\end{frame}

\begin{frame}[fragile]{(4) RUNS = 500}
\lstinputlisting[basicstyle=\ttfamily\fontsize{8pt}{9.6pt}\selectfont,firstline=54,lastline=54]{../src/success_rate_04_the_budget.py}
\end{frame}

\resultframe{The budget}{../figures/success_rate_04_the_budget.png}{%
    Population 10→30: success 80.0\%→99.4\%, cost 100→300 evaluations, yield falls 8.00→3.31 successes per 1000 evaluations; blind draws at equal budget reach 76.2\% and 98.7\% — the GA's margin is +3.8 and +0.7 points}

\section{Next}

\begin{frame}{Next}
On a wide target the GA barely beats blind search at equal budget --- a property of the problem, not a defect of the algorithm.

\vspace{0.8em}
\textbf{Lesson 07 --- Parameter Tuning}

The same instruments, now on the three knobs: crossover, mutation, population. Every knob is secretly a budget knob until evaluations are counted.
\end{frame}
\end{document}
